\subsection{All nonresonant coefficient primes for ordinary RH}
\label{sec:rh-selmer-nonresonant}

For ordinary RH one fixed quadratic field removes the exceptional-set
restriction from the nonresonant estimate.  Keep the actual fields
\[
 E=\mathbf Q(\mu_{2^a})^+(\sqrt{-q},\sqrt{-d}),\qquad
 N=2^aqd,\qquad \Delta=\operatorname{Gal}(E/\mathbf Q),
                                                               \tag{RN1}
\]
where $a\ge2$ and $q,d>3$ are the distinct original primes,
both congruent to $3$ modulo $4$.  For an odd prime $\ell$ and
$e\ge1$, let
\[
 S_{\ell,e}=\ker\!\left(
 H^1(E,\mathbf Q_\ell/\mathbf Z_\ell(-e))
   \longrightarrow\prod_{v<\infty}
                 H^1(E_v,\mathbf Q_\ell/\mathbf Z_\ell(-e))
                         \right),\qquad
 C_e=\frac{1+(-1)^ej}{2},                                    \tag{RN2}
\]
where $j$ is complex conjugation, with its original action on
cohomology and on the coefficient twist.  These are the same
strict Selmer groups and integral parity projections as in
Proposition~\ref{shared-rh-integral-product}.
Indeed, at a finite place away from $\ell$, the unramified
first cohomology is zero: its description as the quotient
by $\operatorname{Frob}_v-1$ vanishes because the nonzero
scalar $(Nv)^{-e}-1$ acts surjectively on
$\mathbf Q_\ell/\mathbf Z_\ell(-e)$.

\begin{theorem}[Uniform bound away from resonance]
\label{thm:rh-selmer-all-nonresonant}
There are absolute constants $C,A>0$ such that for every
field (RN1), every $e\ge1$, and every finite set $L$ of odd
coefficient primes with $\ell-1\nmid2e$, one has
\[
 \sum_{\ell\in L}(\log\ell)
           \operatorname{length}_{\mathbf Z_\ell}(C_eS_{\ell,e})
              \le C N[N(e+1)]^A.                            \tag{RN3}
\]
The constants are independent of $a,q,d,e,L$.  In particular,
for $e=1$ this bounds the sum over every odd coefficient
prime except $3$, with no condition on the character's
$2$-power conductor.
\end{theorem}

\begin{proof}
Use $K=\mathbf Q(\sqrt{-3})$, which has discriminant $-3$
and class number one.  The latter fact can be seen directly:
for $\omega=(-1+\sqrt{-3})/2$, every complex number differs
from an element of $\mathbf Z[\omega]$ by $x+y\omega$ with
$|x|,|y|\le1/2$.  Its squared absolute value is
$x^2-xy+y^2\le3/4<1$.  The norm is therefore Euclidean.

Let $\psi$ be any character of $\Delta$ with
$\psi(-1)=(-1)^e$.  Every such character has power-of-two
order.  At an odd prime, its local conductor component is
quadratic or trivial, since the cyclotomic factor of (RN1)
is ramified only at $2$ and the other factors are quadratic.
Thus the character-order, local-component and parity
conditions in (SP18a) all hold.  The nonresonance condition
forces $\ell\ge5$, so $\ell\nmid d_Kh_K$.
Also $3\nmid N$.

The field $\mathbf Q(\mu_\ell)F_{\psi^{(\ell)}}$ is
unramified at $3$: its two factors have conductors prime
to $3$.  Since $K$ ramifies at $3$, it is not contained
in that compositum.  This verifies the field-exclusion
condition in (SP18a), without imposing
$\ell\notin P$.  Every original good-data hypothesis
is now satisfied for this one fixed $K$
\cite{rh-selmer-nonresonant-source}.

Choose the two auxiliary primes $13$ and $17$.  They do
not divide $30Nd_K$, since $q,d\equiv3\pmod4$ whereas
$13,17\equiv1\pmod4$.  Corollary~\ref{cor:support-spectrum-orbit}
therefore gives its two-prime estimate for every character
orbit and every coefficient prime in $L$, with constants
depending only on $\mathbf Q(\sqrt{-3})$.
Group the characters by their kernels in
$(\mathbf Z/N\mathbf Z)^\times$ and apply (SP24).
The orbit sum is at most $N$, so the total character
length is at most
$2C_KN[N(e+1)]^{A_K}(\log13+\log17)$.
Finally the integral Fourier identity in
Proposition~\ref{shared-rh-integral-product} identifies
that character sum with the length on the left of (RN3),
including its coefficient ramification normalization.
Absorb the two fixed logarithms into $C$.

For $e=1$, an odd prime satisfies $\ell-1\mid2$ exactly
when it is $3$.  The final assertion follows.
\end{proof}

\paragraph{Reduction to the resonant RH estimates.}
Theorem~\ref{thm:rh-selmer-all-nonresonant} also bounds
the nonresonant part of Case~II inside the original
exceptional set, and all of Case~V at $\ell\ge5$.
The remaining coefficient primes satisfy
\[
                        \ell-1\mid2e,
               \qquad \ell\le2e+1.                           \tag{RN4}
\]
Thus for twists $e<t$ all remaining odd coefficient
primes are at most $2t-1$, independently of the growing
cyclotomic layer.  The original size conditions give
$q,d>2t+1$ and $q,d>\max P$, where $P$ is the original
finite exceptional set\cite{rh-selmer-nonresonant-source}.
Hence every remaining resonant
prime is prime to $N$.  At such a prime the local conductor
component of every character is trivial.  Every nontrivial
character of the $2$-group $\Delta$ has even order.
Consequently the nontrivial characters with $e\ge2$
are exactly the remaining Case~II characters; the trivial
character is covered by the existing fixed-field contribution.

Here is one precise sufficient frontier for the residual
Selmer estimates.  For $e\ge2$ put
\[
       T_{\ell,e}=C_e(1-A_\Delta)S_{\ell,e},\qquad
       A_\Delta=|\Delta|^{-1}\sum_{g\in\Delta}g.              \tag{RN5}
\]
For $e=1$ the only remaining prime is $3$.  The original
threshold is
\[
 R_3=16\cdot3^2\cdot4^{v_2(3^2-1)}=9216,
       \qquad \max\{c:2^c<R_3\}=13.
\]
Set
\[
 E_{3,\mathrm{low}}=
       \mathbf Q(\mu_{2^{\min(a,13)}})^+(\sqrt{-q},\sqrt{-d}),
 \quad H_3=\operatorname{Gal}(E/E_{3,\mathrm{low}}),
 \quad T_{3,1}=C_1(1-A_{H_3})S_{3,1}.                         \tag{RN6}
\]
The field in (RN6) retains the original $q,d$; it is not
one absolute field as those primes vary.

\emph{Combined resonant estimate.}  The desired estimate is that there is a constant
$B>0$, independent of the original family and of $e$,
such that in the original twist range, for every odd
prime satisfying (RN4),
\[
             (\log\ell)\,
             \operatorname{length}_{\mathbf Z_\ell}T_{\ell,e}
                         \le[N(e+1)]^B.                      \tag{RN7}
\]
The original Case~II character bound, summed over at most
$N$ characters, and the original Case~V bound at $3$
would imply (RN7).  Conversely, the integral direct-sum
decomposition makes each relevant character length no
larger than the total in (RN7).  Thus a proof of (RN7)
would supply both remaining RH Selmer estimates, after
the already established fixed-field and small-conductor
contributions are included.  All the groups in (RN5)--(RN6)
retain their integral torsion.

For precision, the source requires one exponent $C$ for
all hosts of the fixed target, all $t$ satisfying the
original size conditions, all $0<e<t$, and all coefficient
primes and characters in the stated cases.  Its individual
bound is $(\ell Ne)^C$\cite{rh-selmer-nonresonant-source}.
Since $N\ge2$ and $\ell\le2e+1$, summing at most $N$
characters gives
\[
 N(\ell Ne)^C\le[N(e+1)]^{2C+1}.
\]
In the other direction $N(e+1)\le\ell Ne$, so (RN7)
gives the original individual bound with exponent $B$.
These inequalities preserve the required uniform quantifiers.

Corollary~\ref{cor:rh-selmer-exponent-one} below supplies (RN7)
for $e=1$, and the high-dyadic-conductor part of (RN5) at three
for every $e\ge2$. Thus the remaining frontier is (RN7) for
$e\ge2$, with only the characters of dyadic conductor at most
$64$ retained at $\ell=3$. The other terms already have a
uniform polynomial bound.
